\documentclass[12pt,a4paper]{report}
\usepackage[T1]{fontenc}
\usepackage[french]{babel}
\usepackage{geometry}
\usepackage{graphicx}
\usepackage{fancyhdr}
\usepackage{listings}
\usepackage{xcolor}
\usepackage{hyperref}
\usepackage{array}
\usepackage{tabularx}
\usepackage{multicol}
\usepackage{fancybox}
\usepackage{tocloft}
\usepackage{float}
\usepackage{comment}
\usepackage{amsmath}

% Configuration de la géométrie
\geometry{left=2.5cm, right=2.5cm, top=2.5cm, bottom=2.5cm}

% Configuration des en-têtes et pieds de page
\pagestyle{fancy}
\fancyhf{}
\rhead{\textit{Projet DevOps - Cloud Infrastructure}}
\lhead{\textit{Rapport Académique}}
\cfoot{\thepage}
\renewcommand{\headrulewidth}{0.4pt}
\renewcommand{\footrulewidth}{0.4pt}

% Configuration des listings
\lstset{
    basicstyle=\ttfamily\small,
    keywordstyle=\color{blue},
    commentstyle=\color{gray},
    stringstyle=\color{red},
    breaklines=true,
    showstringspaces=false,
    frame=tb,
    rulecolor=\color{black},
    tabsize=4
}

% Titre et auteurs
\title{
    \textbf{RAPPORT DE PROJET DEVOPS} \\
    \textbf{Infrastructure Cloud CI/CD et Déploiement Kubernetes} \\[1cm]
    {\large Automating Cloud Infrastructure with Docker, Terraform, Ansible and Kubernetes}
}

\author{
    \textbf{Auteurs:} \\
    LEGNIOUI SIFEDDINE \\
    LOUAT OUSSAMA \\
    AYOUB AFRAOUI \\
    KAWTAR ELBEJJAJI \\[1cm]
    \textit{MEIA - Master Expertises en Informatique Avancée}
}

\date{\today}

\begin{document}

% Page de titre
\begin{titlepage}
    \centering
    \vspace*{2cm}
    {\Large \textbf{UNIVERSITÉ SIDI MOHAMMED BEN ABDELLAH}} \\[0.5cm]
    {\Large \textbf{FACULTÉ DES SCIENCES}} \\[0.5cm]
    {\Large Département d'Informatique} \\[3cm]
    
    \rule{\linewidth}{1.5pt} \\[0.5cm]
    {\LARGE \textbf{RAPPORT DE PROJET}} \\[0.2cm]
    {\LARGE \textbf{INFRASTRUCTURE DEVOPS}} \\
    {\LARGE \textbf{ET CLOUD COMPUTING}} \\[0.5cm]
    \rule{\linewidth}{1.5pt} \\[2cm]
    
    {\Large \textbf{Automating Cloud Infrastructure with Docker,}} \\
    {\Large \textbf{Terraform, Ansible and Kubernetes}} \\[3cm]
    
    {\large \textbf{Auteurs:}} \\[0.3cm]
    {\large LEGNIOUI SIFEDDINE} \\
    {\large LOUAT OUSSAMA} \\
    {\large AYOUB AFRAOUI} \\
    {\large KAWTAR ELBEJJAJI} \\[2cm]
    
    {\large \textbf{Program:}} Master Expertises en Informatique Avancée (MEIA) \\[1cm]
    
    {\large \textbf{Date:}} \today \\[1cm]
    
    {\large \textbf{Année Académique:} 2025-2026}
    
    \vfill
    
    {\large Fès, Maroc}
\end{titlepage}

% Table des matières
\tableofcontents
\newpage

% Résumé Exécutif
\chapter*{Résumé Exécutif}
\addcontentsline{toc}{chapter}{Résumé Exécutif}

Ce rapport présente une solution complète d'infrastructure cloud et DevOps, combinant les meilleures pratiques modernes en matière de conteneurisation, orchestration et infrastructure-as-code.

Le projet implémente une pipeline CI/CD entièrement automatisée, permettant le déploiement rapide et fiable d'applications web, tout en assurant la scalabilité, la haute disponibilité et la maintenabilité de l'infrastructure.

\textbf{Principaux objectifs réalisés:}
\begin{enumerate}
    \item Développement d'une application web Python moderne
    \item Conteneurisation avec Docker pour portabilité
    \item Infrastructure-as-Code avec Terraform
    \item Automatisation avec Ansible pour configuration/déploiement
    \item Orchestration Kubernetes pour haute disponibilité
    \item Pipeline CI/CD intégré pour déploiement continu
    \item Documentation complète et reproductibilité
\end{enumerate}

\textbf{Technologies clés:}
\begin{itemize}
    \item \textbf{Application:} Python Flask
    \item \textbf{Conteneurisation:} Docker & Docker Compose
    \item \textbf{Infrastructure:} Terraform (IaaC)
    \item \textbf{Configuration:} Ansible
    \item \textbf{Orchestration:} Kubernetes
    \item \textbf{CI/CD:} GitHub Actions / Pipeline
\end{itemize}

\newpage

% CHAPITRE 1: Introduction
\chapter{Introduction}

\section{Contexte du Projet}

L'infrastructure informatique moderne exige une approche holistique combinant plusieurs disciplines: development, operations, et security. Le paradigme DevOps émerge comme solution pour éliminer les silos organisationnels et accélérer le time-to-market tout en assurant la qualité et la sécurité.

Ce projet s'inscrit dans cette démarche en implémentant une approche complète d'automatisation et d'orchestration sur cloud infrastructure.

\section{Objectifs du Projet}

Les objectifs principaux de ce projet sont:

\begin{itemize}
    \item Concevoir et implémenter une application web scalable et maintenable
    \item Mettre en place une infrastructure cloud automatisée et reproductible
    \item Implémenter une pipeline CI/CD pour automatiser les tests et déploiements
    \item Assurer la haute disponibilité et la résilience de l'application
    \item Documenter toutes les étapes pour faciliter la reproduction et la maintenance
    \item Démontrer les meilleures pratiques DevOps et Cloud Native
\end{itemize}

\section{Périmètre de la Solution}

Cette solution couvre:

\begin{enumerate}
    \item \textbf{Application}: Application Flask Python simple mais fonctionnelle
    \item \textbf{Conteneurisation}: Dockerfile pour packaging de l'application
    \item \textbf{Infrastructure as Code}: Scripts Terraform pour provision cloud
    \item \textbf{Configuration Management}: Playbooks Ansible pour configuration serveurs
    \item \textbf{Orchestration}: Manifests Kubernetes YAML pour déploiement
    \item \textbf{Pipeline CI/CD}: Configuration au niveau VCS pour automatisation
    \item \textbf{Documentation}: Guide complet de reproduction et maintenance
\end{enumerate}

\newpage

% CHAPITRE 2: Architecture et Design
\chapter{Architecture et Design de la Solution}

\section{Vue d'ensemble de l'Architecture}

L'architecture proposée suit une approche microservices avec orchestration Kubernetes, permettant scalabilité horizontale et haute disponibilité.

\subsection{Composants Principaux}

\begin{itemize}
    \item \textbf{Couche Application}: Services applicatifs dockerisés
    \item \textbf{Couche Orchestration}: Kubernetes pour gestion conteneurs
    \item \textbf{Couche Infrastructure}: Cloud provider géré par Terraform
    \item \textbf{Couche Automation}: Ansible pour configuration initiale
    \item \textbf{Couche CI/CD}: Pipeline automatisée pour tests et déploiement
\end{itemize}

\section{Flux de Déploiement}

\begin{enumerate}
    \item Le développeur commit du code sur le repository Git
    \item La pipeline CI/CD est déclenchée automatiquement
    \item Tests unitaires et d'intégration sont exécutés
    \item L'image Docker est construite et pushée vers registry
    \item Kubernetes applique les nouveaux manifests
    \item L'application est mise à jour en production
\end{enumerate}

\section{Choix Technologiques Justifiés}

\subsection{Docker pour Conteneurisation}

Docker fournit:
\begin{itemize}
    \item Isolation des applications et dépendances
    \item Portabilité garantie entre environnements
    \item Réduction overhead vs machines virtuelles
    \item Standard industrie pour conteneurisation
    \item Écosystème riche (Docker Hub, Docker Compose)
\end{itemize}

\subsection{Terraform pour Infrastructure as Code}

Terraform offre:
\begin{itemize}
    \item Infrastructure déclarative et versionnable
    \item Multi-cloud capability (AWS, Azure, GCP)
    \item State management pour tracking changements
    \item Reproducibilité garantie
    \item Facilite destruction et recréation infrastructure
\end{itemize}

\subsection{Ansible pour Configuration Management}

Ansible fournit:
\begin{itemize}
    \item Configuration agentless (SSH uniquement)
    \item Syntax simple et lisible (YAML)
    \item Idempotence pour opérations sûres
    \item Large libraire de modules prêts à l'emploi
    \item Provisioning et gestion configuration unifiée
\end{itemize}

\subsection{Kubernetes pour Orchestration}

Kubernetes apporte:
\begin{itemize}
    \item Orchestration automatique de conteneurs
    \item Haute disponibilité avec réplication
    \item Scalabilité horizontale automatique
    \item Load balancing intégré
    \item Gestion des secrets et configurations
    \item Standard industrie pour conteneurs en production
\end{itemize}

\newpage

% CHAPITRE 3: Application
\chapter{Application Web}

\section{Vue d'ensemble de l'Application}

L'application est une web application Flask Python simple, servant comme cas d'usage démontrant les capacités de l'infrastructure DevOps.

\section{Source Code}

\subsection{Fichier Principal - app.py}

\lstset{language=python}
\begin{lstlisting}
from flask import Flask, render_template
import os

app = Flask(__name__)

@app.route('/')
def home():
    return render_template('index.html')

@app.route('/health')
def health():
    return {'status': 'healthy'}, 200

@app.route('/version')
def version():
    return {'version': '1.0.0'}, 200

if __name__ == '__main__':
    app.run(host='0.0.0.0', port=5000, debug=False)
\end{lstlisting}

\subsection{Dépendances - requirements.txt}

\begin{lstlisting}
Flask==2.3.0
Werkzeug==2.3.0
\end{lstlisting}

\subsection{Template HTML - index.html}

%\lstset{language=html}
\begin{lstlisting}
<!DOCTYPE html>
<html>
<head>
    <title>DevOps Project</title>
    <link rel="stylesheet" href="/static/style.css">
</head>
<body>
    <div class="container">
        <h1>DevOps Cloud Infrastructure Project</h1>
        <p>Application déployée avec succès!</p>
        <ul>
            <li>Docker Containerization</li>
            <li>Kubernetes Orchestration</li>
            <li>Terraform IaC</li>
            <li>Ansible Configuration</li>
        </ul>
    </div>
    <script src="/static/script.js"></script>
</body>
</html>
\end{lstlisting}

\subsection{Styles CSS - style.css}

%\lstset{language=css}
\begin{lstlisting}
body {
    font-family: Arial, sans-serif;
    background-color: #f4f4f4;
    margin: 0;
    padding: 20px;
}

.container {
    max-width: 800px;
    margin: 0 auto;
    background-color: white;
    padding: 30px;
    border-radius: 8px;
    box-shadow: 0 2px 10px rgba(0,0,0,0.1);
}

h1 {
    color: #333;
    border-bottom: 3px solid #007bff;
    padding-bottom: 10px;
}

ul {
    list-style-type: none;
    padding: 0;
}

li {
    background-color: #e7f3ff;
    margin: 10px 0;
    padding: 10px;
    border-left: 4px solid #007bff;
}
\end{lstlisting}

\subsection{Interactivité JavaScript - script.js}

%\lstset{language=javascript}
\begin{lstlisting}
document.addEventListener('DOMContentLoaded', function() {
    console.log('Application loaded successfully');
    
    // Appel health check
    fetch('/health')
        .then(response => response.json())
        .then(data => console.log('Health check:', data));
});
\end{lstlisting}

\section{Spécifications Techniques}

\begin{table}[H]
\centering
\begin{tabularx}{\linewidth}{|l|X|}
\hline
\textbf{Paramètre} & \textbf{Valeur} \\
\hline
Langage & Python 3.9+ \\
Framework & Flask 2.3.0 \\
Port Application & 5000 \\
Endpoints & /, /health, /version \\
Dépendances Externes & Minimales \\
\hline
\end{tabularx}
\end{table}

\newpage

% CHAPITRE 4: Conteneurisation Docker
\chapter{Conteneurisation - Docker et Dockerfiles}

\section{Concept et Avantages}

Docker est une plateforme de conteneurisation qui encapsule l'application et toutes ses dépendances dans une image repérable et portable.

\subsection{Avantages de Docker}

\begin{itemize}
    \item \textbf{Portabilité}: "Build once, run anywhere"
    \item \textbf{Isolation}: Applications isolées les unes des autres
    \item \textbf{Efficacité}: Moins de ressources que VMs
    \item \textbf{Rapidité}: Démarrage en millisecondes
    \item \textbf{Scalabilité}: Facile de créer/détruire conteneurs
\end{itemize}

\section{Dockerfile}

Le Dockerfile définit l'image Docker qui sera utilisée pour exécuter l'application.

%\lstset{language=docker)
\begin{lstlisting}
# Utiliser image officielle Python
FROM python:3.9-slim

# Définir répertoire de travail
WORKDIR /app

# Copier requirements et installer dépendances
COPY requirements.txt .
RUN pip install --no-cache-dir -r requirements.txt

# Copier code source
COPY app.py .
COPY templates/ templates/
COPY static/ static/

# Exposer port
EXPOSE 5000

# Healthcheck
HEALTHCHECK --interval=30s --timeout=3s --start-period=5s --retries=3 \
    CMD python -c "import urllib.request; \
    urllib.request.urlopen('http://localhost:5000/health').read()"

# Commande de démarrage
CMD ["python", "app.py"]
\end{lstlisting}

\section{Bonnes Pratiques Docker}

\subsection{Optimisation des Images}

\begin{enumerate}
    \item \textbf{Base image minimal}: Utiliser `python:3.9-slim` au lieu de `python:3.9`
    \item \textbf{Multi-stage builds}: Réduire taille image finale
    \item \textbf{Layer caching}: Ordonner instructions pour optimiser cache
    \item \textbf{Cleanup}: Nettoyer après installation depandances
\end{enumerate}

\subsection{Sécurité Docker}

\begin{enumerate}
    \item \textbf{User non-root}: Ne pas exécuter application en tant que root
    \item \textbf{Scans vulnérabilités}: Scanner images régulièrement
    \item \textbf{Update base images}: Garder images à jour
    \item \textbf{Secrets management}: Ne pas hardcoder secrets
\end{enumerate}

\section{Docker Compose}

Docker Compose permet de définir et exécuter multi-conteneurs applications localement:

\lstset{language=yaml}
\begin{lstlisting}
version: '3.8'

services:
  app:
    build: .
    ports:
      - "5000:5000"
    environment:
      - FLASK_ENV=production
    volumes:
      - ./app:/app
    healthcheck:
      test: ["CMD", "curl", "-f", "http://localhost:5000/health"]
      interval: 30s
      timeout: 10s
      retries: 3
      start_period: 5s
\end{lstlisting}

\section{Commandes Docker Essentielles}

\begin{table}[H]
\centering
\begin{tabularx}{\linewidth}{|l|X|}
\hline
\textbf{Commande} & \textbf{Description} \\
\hline
docker build -t app:1.0 . & Construire image \\
docker run -p 5000:5000 app:1.0 & Exécuter conteneur \\
docker push registry/app:1.0 & Pusher image vers registry \\
docker pull app:1.0 & Télécharger image \\
docker ps & Lister conteneurs actifs \\
docker logs <container\_id> & Voir logs conteneur \\
docker exec -it <container\_id> bash & Accéder shell conteneur \\
\hline
\end{tabularx}
\end{table}

\newpage

% CHAPITRE 5: Infrastructure as Code - Terraform
\chapter{Infrastructure as Code - Terraform}

\section{Introduction à Terraform}

Terraform est un outil Infrastructure as Code (IaC) permettant de:
\begin{itemize}
    \item Définir infrastructure en code déclaratif
    \item Versionner et contrôler infrastructure
    \item Déployer infrastructure reproduite sur multiples environnements
    \item Faciliter destruction et recréation clean
\end{itemize}

\section{Structure Terraform}

\subsection{providers.tf}

Définit les providers utilisés (AWS, Azure, GCP, etc.):

\lstset{language=hcl}
\begin{lstlisting}
terraform {
  required_version = ">= 1.0"
  
  required_providers {
    aws = {
      source  = "hashicorp/aws"
      version = "~> 5.0"
    }
  }
}

provider "aws" {
  region = var.aws_region

  default_tags {
    tags = {
      Project     = "devops-project"
      Environment = var.environment
      ManagedBy   = "Terraform"
    }
  }
}
\end{lstlisting}

\subsection{variables.tf}

Définit les variables d'entrée:

\lstset{language=hcl}
\begin{lstlisting}
variable "aws_region" {
  description = "AWS region"
  type        = string
  default     = "us-east-1"
}

variable "environment" {
  description = "Environment name"
  type        = string
  default     = "production"
}

variable "app_name" {
  description = "Application name"
  type        = string
  default     = "devops-app"
}

variable "instance_type" {
  description = "EC2 instance type"
  type        = string
  default     = "t3.medium"
}

variable "cluster_size" {
  description = "Kubernetes cluster size"
  type        = number
  default     = 3
}
\end{lstlisting}

\subsection{main.tf}

Ressources principales:

\lstset{language=hcl}
\begin{lstlisting}
# VPC
resource "aws_vpc" "main" {
  cidr_block           = "10.0.0.0/16"
  enable_dns_hostnames = true
  enable_dns_support   = true

  tags = {
    Name = "${var.app_name}-vpc"
  }
}

# Subnets publics
resource "aws_subnet" "public" {
  count                   = 2
  vpc_id                  = aws_vpc.main.id
  cidr_block              = "10.0.${count.index + 1}.0/24"
  availability_zone       = data.aws_availability_zones.available.names[count.index]
  map_public_ip_on_launch = true

  tags = {
    Name = "${var.app_name}-public-subnet-${count.index + 1}"
  }
}

# Internet Gateway
resource "aws_internet_gateway" "main" {
  vpc_id = aws_vpc.main.id

  tags = {
    Name = "${var.app_name}-igw"
  }
}

# EKS Cluster
resource "aws_eks_cluster" "main" {
  name     = "${var.app_name}-cluster"
  role_arn = aws_iam_role.eks_role.arn
  version  = "1.27"

  vpc_config {
    subnet_ids = aws_subnet.public[*].id
  }

  depends_on = [
    aws_iam_role_policy_attachment.eks_cluster_policy
  ]
}

# EKS Node Group
resource "aws_eks_node_group" "main" {
  cluster_name    = aws_eks_cluster.main.name
  node_group_name = "${var.app_name}-node-group"
  node_role_arn   = aws_iam_role.node_role.arn
  subnet_ids      = aws_subnet.public[*].id

  scaling_config {
    desired_size = var.cluster_size
    max_size     = var.cluster_size + 1
    min_size     = var.cluster_size
  }

  instance_types = [var.instance_type]

  depends_on = [
    aws_iam_role_policy_attachment.eks_node_policy
  ]
}
\end{lstlisting}

\subsection{outputs.tf}

Exporte les informations importantes:

\lstset{language=hcl}
\begin{lstlisting}
output "eks_cluster_endpoint" {
  description = "EKS cluster API endpoint"
  value       = aws_eks_cluster.main.endpoint
}

output "eks_cluster_name" {
  description = "EKS cluster name"
  value       = aws_eks_cluster.main.name
}

output "kubeconfig" {
  description = "kubectl config"
  value       = "aws eks update-kubeconfig --region ${var.aws_region} --name ${aws_eks_cluster.main.name}"
  sensitive   = false
}
\end{lstlisting}

\section{Flux Terraform Typique}

\begin{enumerate}
    \item \texttt{terraform init}: Initialiser le répertoire
    \item \texttt{terraform plan}: Afficher changements proposés
    \item \texttt{terraform apply}: Appliquer changements
    \item \texttt{terraform destroy}: Détruire les ressources
\end{enumerate}

\section{Bonnes Pratiques Terraform}

\begin{itemize}
    \item \textbf{State management}: Stocker état sur backend remote (S3, Terraform Cloud)
    \item \textbf{Variables}: Utiliser fichiers .tfvars pour valeurs sensibles
    \item \textbf{Modules}: Organiser code en modules réutilisables
    \item \textbf{Linting}: Valider syntaxe avec \texttt{terraform fmt} et \texttt{terraform validate}
    \item \textbf{Plan before apply}: Toujours vérifier avec \texttt{terraform plan}
\end{itemize}

\newpage

% CHAPITRE 6: Configuration Management - Ansible
\chapter{Configuration et Automatisation - Ansible}

\section{Vue d'ensemble Ansible}

Ansible est un outil de configuration management et orchestration permettant:
\begin{itemize}
    \item Automatiser déploiement et configuration
    \item Gérer infrastructures complexes à grande échelle
    \item Rester agentless (SSH uniquement requis)
    \item Écrire playbooks lisibles en YAML
\end{itemize}

\section{Inventaire}

Le fichier inventory.ini définit les hôtes gérés:

%\lstset{language=ini)
\begin{lstlisting}
[all]
app_server_1 ansible_host=10.0.1.100 ansible_user=ubuntu
app_server_2 ansible_host=10.0.2.100 ansible_user=ubuntu
db_server ansible_host=10.0.3.100 ansible_user=ubuntu

[app]
app_server_1
app_server_2

[database]
db_server

[all:vars]
ansible_python_interpreter=/usr/bin/python3
ansible_ssh_private_key_file=~/.ssh/id_rsa

[app:vars]
docker_registry=registry.example.com
app_version=1.0.0
\end{lstlisting}

\section{Playbook Principal}

\lstset{language=yaml}
\begin{lstlisting}
---
- name: Configure DevOps Infrastructure
  hosts: all
  become: yes
  
  tasks:
    - name: Update system packages
      apt:
        update_cache: yes
        cache_valid_time: 3600
      when: ansible_os_family == "Debian"

    - name: Install required packages
      package:
        name:
          - curl
          - wget
          - git
          - htop
          - net-tools
        state: present

    - name: Install Docker
      block:
        - name: Add Docker GPG key
          apt_key:
            url: https://download.docker.com/linux/ubuntu/gpg
            state: present

        - name: Add Docker repository
          apt_repository:
            repo: "deb [arch=amd64] https://download.docker.com/linux/ubuntu {{ ansible_distribution_release }} stable"
            state: present

        - name: Install Docker
          apt:
            name: docker-ce
            state: present

        - name: Start Docker service
          service:
            name: docker
            state: started
            enabled: yes

    - name: Install Kubernetes tools
      block:
        - name: Install kubectl
          shell: |
            curl -LO "https://dl.k8s.io/release/$(curl -L -s https://dl.k8s.io/release/stable.txt)/bin/linux/amd64/kubectl"
            chmod +x kubectl
            mv kubectl /usr/local/bin/

        - name: Install helm
          shell: |
            curl https://raw.githubusercontent.com/helm/helm/main/scripts/get-helm-3 | bash

    - name: Deploy application
      block:
        - name: Clone application repository
          git:
            repo: "{{ git_repo }}"
            dest: /opt/app
            version: main

        - name: Build Docker image
          shell: docker build -t app:latest /opt/app

        - name: Run application container
          docker_container:
            name: devops_app
            image: app:latest
            ports:
              - "5000:5000"
            state: started
            restart_policy: always
\end{lstlisting}

\section{Modules Ansible Essentiels}

\begin{table}[H]
\centering
\begin{tabularx}{\linewidth}{|l|X|}
\hline
\textbf{Module} & \textbf{Utilisation} \\
\hline
apt & Gestion paquets Debian/Ubuntu \\
yum & Gestion paquets RedHat/CentOS \\
service & Gestion services système \\
copy & Copier fichiers vers serveurs \\
template & Déployer templates avec variables \\
shell & Exécuter commands shell \\
docker\_container & Gérer conteneurs Docker \\
kubernetes.core.k8s & Gérer ressources Kubernetes \\
\hline
\end{tabularx}
\end{table}

\section{Exécution des Playbooks}

\begin{lstlisting}
# Exécuter playbook entier
ansible-playbook playbook.yml

# Exécuter playbook sur hôtes spécifiques
ansible-playbook playbook.yml -i inventory.ini -l app

# Mode dry-run (sans appliquer changements)
ansible-playbook playbook.yml --check

# Verbose mode pour debugging
ansible-playbook playbook.yml -vvv

# Utiliser vault pour données sensibles
ansible-playbook playbook.yml --ask-vault-pass
\end{lstlisting}

\newpage

% CHAPITRE 7: Orchestration Kubernetes
\chapter{Orchestration - Kubernetes et Manifests}

\section{Vue d'ensemble Kubernetes}

Kubernetes est un système open-source d'orchestration de conteneurs permettant:
\begin{itemize}
    \item Automatiser déploiement et gestion conteneurs
    \item Scalabilité horizontale automatique
    \item Haute disponibilité et résilience
    \item Load balancing et découverte services
    \item Gestion ressources et scheduling intelligent
\end{itemize}

\section{Concepts Clés Kubernetes}

\subsection{Pod}
Unité de déploiement minimale, encapsulant un ou plusieurs conteneurs.

\subsection{Deployment}
Gère réplication et updates de Pods de façon déclarative.

\subsection{Service}
Expose Pods de manière stable via IP et DNS.

\subsection{ConfigMap et Secret}
Gèrent respectivement configurations et données sensibles.

\subsection{Namespace}
Fournit isolation logique entre ressources.

\section{Deployment Manifest}

\lstset{language=yaml}
\begin{lstlisting}
apiVersion: apps/v1
kind: Deployment
metadata:
  name: devops-app
  namespace: production
  labels:
    app: devops-app
    version: v1
spec:
  replicas: 3
  
  strategy:
    type: RollingUpdate
    rollingUpdate:
      maxSurge: 1
      maxUnavailable: 0
  
  selector:
    matchLabels:
      app: devops-app
  
  template:
    metadata:
      labels:
        app: devops-app
        version: v1
      annotations:
        prometheus.io/scrape: "true"
        prometheus.io/port: "5000"
    
    spec:
      serviceAccountName: devops-app
      
      securityContext:
        runAsNonRoot: true
        runAsUser: 1000
        fsGroup: 1000
      
      containers:
      - name: app
        image: devops-app:latest
        imagePullPolicy: Always
        
        ports:
        - name: http
          containerPort: 5000
          protocol: TCP
        
        env:
        - name: FLASK_ENV
          value: "production"
        - name: LOG_LEVEL
          value: "INFO"
        
        envFrom:
        - configMapRef:
            name: app-config
        
        resources:
          requests:
            cpu: 100m
            memory: 128Mi
          limits:
            cpu: 500m
            memory: 512Mi
        
        livenessProbe:
          httpGet:
            path: /health
            port: http
          initialDelaySeconds: 30
          periodSeconds: 10
          timeoutSeconds: 3
          failureThreshold: 3
        
        readinessProbe:
          httpGet:
            path: /health
            port: http
          initialDelaySeconds: 5
          periodSeconds: 5
          timeoutSeconds: 1
          failureThreshold: 2
        
        volumeMounts:
        - name: config
          mountPath: /app/config
          readOnly: true
      
      volumes:
      - name: config
        configMap:
          name: app-config
      
      affinity:
        podAntiAffinity:
          preferredDuringSchedulingIgnoredDuringExecution:
          - weight: 100
            podAffinityTerm:
              labelSelector:
                matchExpressions:
                - key: app
                  operator: In
                  values:
                  - devops-app
              topologyKey: kubernetes.io/hostname
      
      terminationGracePeriodSeconds: 30
\end{lstlisting}

\section{Service Manifest}

\lstset{language=yaml}
\begin{lstlisting}
apiVersion: v1
kind: Service
metadata:
  name: devops-app
  namespace: production
  labels:
    app: devops-app
spec:
  type: LoadBalancer
  selector:
    app: devops-app
  ports:
  - name: http
    protocol: TCP
    port: 80
    targetPort: 5000
  sessionAffinity: None
  
  # Alternative: ClusterIP pour interne seulement
  # type: ClusterIP
  # selector:
  #   app: devops-app
  # ports:
  # - name: http
  #   protocol: TCP
  #   port: 5000
  #   targetPort: 5000
\end{lstlisting}

\section{ConfigMap et Secret}

\subsection{ConfigMap}

\lstset{language=yaml}
\begin{lstlisting}
apiVersion: v1
kind: ConfigMap
metadata:
  name: app-config
  namespace: production
data:
  app.properties: |
    environment=production
    debug=false
    log_level=INFO
  database_url: "postgresql://db:5432/appdb"
\end{lstlisting}

\subsection{Secret}

\lstset{language=yaml}
\begin{lstlisting}
apiVersion: v1
kind: Secret
metadata:
  name: app-secrets
  namespace: production
type: Opaque
stringData:
  db_username: "appuser"
  db_password: "secure_password_here"
  api_key: "secret_api_key"
\end{lstlisting}

\section{Ingress pour Routage HTTP}

\lstset{language=yaml}
\begin{lstlisting}
apiVersion: networking.k8s.io/v1
kind: Ingress
metadata:
  name: devops-app-ingress
  namespace: production
  annotations:
    cert-manager.io/cluster-issuer: "letsencrypt-prod"
spec:
  ingressClassName: nginx
  tls:
  - hosts:
    - app.example.com
    secretName: app-tls-cert
  rules:
  - host: app.example.com
    http:
      paths:
      - path: /
        pathType: Prefix
        backend:
          service:
            name: devops-app
            port:
              number: 80
\end{lstlisting}

\section{Commandes kubectl Essentielles}

\begin{table}[H]
\centering
\begin{tabularx}{\linewidth}{|l|X|}
\hline
\textbf{Commande} & \textbf{Description} \\
\hline
kubectl apply -f deployment.yaml & Créer/mettre à jour ressources \\
kubectl get pods & Lister les pods \\
kubectl describe pod <pod\_name> & Détails d'un pod \\
kubectl logs <pod\_name> & Voir logs d'un pod \\
kubectl exec -it <pod> bash & Shell dans un pod \\
kubectl port-forward <pod> 5000:5000 & Tunneler port local \\
kubectl rollout status deployment/app & Statut deployment \\
kubectl scale deployment app --replicas=5 & Scaler replicas \\
kubectl delete -f deployment.yaml & Supprimer ressources \\
\hline
\end{tabularx}
\end{table}

\newpage

% CHAPITRE 8: Pipeline CI/CD
\chapter{Pipeline d'Intégration et Déploiement Continus}

\section{Concept et Architecture CI/CD}

Une pipeline CI/CD automatise:
\begin{enumerate}
    \item \textbf{Continuous Integration}: Tests automatiques à la moindre modification
    \item \textbf{Continuous Delivery}: Prepare releases prêtes à déployer
    \item \textbf{Continuous Deployment}: Déploie automatiquement en production
\end{enumerate}

\section{Flux Typique de Pipeline}

\begin{enumerate}
    \item \textbf{Trigger}: Code pushed au repository
    \item \textbf{Checkout}: Récupérer code source
    \item \textbf{Build}: Compiler/construire application
    \item \textbf{Test}: Exécuter tests unitaires et intégration
    \item \textbf{Code Quality}: Analyser qualité code
    \item \textbf{Build Image}: Créer image Docker
    \item \textbf{Push Registry}: Pusher vers container registry
    \item \textbf{Deploy}: Déployer en staging
    \item \textbf{Integration Tests}: Tests en environnement déployé
    \item \textbf{Approve}: Validation manuelle (optionnel)
    \item \textbf{Production Deploy}: Déployer en production
    \item \textbf{Smoke Tests}: Tests de santé en production
    \item \textbf{Notification}: Notifier équipes
\end{enumerate}

\section{GitHub Actions Configuration}

\lstset{language=yaml}
\begin{lstlisting}
name: CI/CD Pipeline

on:
  push:
    branches: [ main, develop ]
  pull_request:
    branches: [ main ]

env:
  REGISTRY: ghcr.io
  IMAGE_NAME: ${{ github.repository }}

jobs:
  test:
    runs-on: ubuntu-latest
    
    steps:
    - uses: actions/checkout@v3
    
    - name: Set up Python
      uses: actions/setup-python@v4
      with:
        python-version: '3.9'
    
    - name: Install dependencies
      run: |
        pip install -r requirements.txt
        pip install pytest pytest-cov
    
    - name: Run tests
      run: pytest tests/ -v --cov
    
    - name: Upload coverage
      uses: codecov/codecov-action@v3

  build:
    needs: test
    runs-on: ubuntu-latest
    permissions:
      contents: read
      packages: write
    
    if: github.event_name == 'push'
    
    steps:
    - uses: actions/checkout@v3
    
    - name: Set up Docker Buildx
      uses: docker/setup-buildx-action@v2
    
    - name: Log in to Registry
      uses: docker/login-action@v2
      with:
        registry: ${{ env.REGISTRY }}
        username: ${{ github.actor }}
        password: ${{ secrets.GITHUB_TOKEN }}
    
    - name: Extract metadata
      id: meta
      uses: docker/metadata-action@v4
      with:
        images: ${{ env.REGISTRY }}/${{ env.IMAGE_NAME }}
        tags: |
          type=ref,event=branch
          type=sha,prefix={{branch}}-
          type=semver,pattern={{version}}
    
    - name: Build and push Docker image
      uses: docker/build-push-action@v4
      with:
        context: .
        push: true
        tags: ${{ steps.meta.outputs.tags }}
        labels: ${{ steps.meta.outputs.labels }}

  deploy:
    needs: build
    runs-on: ubuntu-latest
    if: github.ref == 'refs/heads/main'
    
    steps:
    - uses: actions/checkout@v3
    
    - name: Deploy to Kubernetes
      env:
        KUBE_CONFIG: ${{ secrets.KUBE_CONFIG }}
      run: |
        mkdir -p $HOME/.kube
        echo "$KUBE_CONFIG" | base64 -d > $HOME/.kube/config
        kubectl apply -f kubernetes/
        kubectl rollout status deployment/devops-app
    
    - name: Run smoke tests
      run: |
        ./scripts/smoke_tests.sh
    
    - name: Notify Slack
      if: always()
      uses: slackapi/slack-github-action@v1
      with:
        webhook-url: ${{ secrets.SLACK_WEBHOOK }}
        payload: |
          {
            "text": "Deployment ${{ job.status }}",
            "blocks": [
              {
                "type": "section",
                "text": {
                  "type": "mrkdwn",
                  "text": "*Deployment Status*: ${{ job.status }}"
                }
              }
            ]
          }
\end{lstlisting}

\section{Bonnes Pratiques CI/CD}

\begin{itemize}
    \item \textbf{Test First}: Écrire tests avant/pendant code
    \item \textbf{Automated}: Tout automatisé autant que possible
    \item \textbf{Fast Feedback}: Résultats rapidement pour itération rapide
    \item \textbf{Stage Environment}: Tester en env similaire à production
    \item \textbf{Rollback Ready}: Capacité de rollback rapide en cas problème
    \item \textbf{Monitoring}: Alertes détaillées sur chaque étape
    \item \textbf{Documentation}: Dépendances et requirements documentés
\end{itemize}

\newpage

% CHAPITRE 9: Gestion des Secrets et Sécurité
\chapter{Gestion des Secrets et Sécurité}

\section{Principes de Sécurité}

\subsection{Ne jamais hardcoder secrets}

Les secrets (passwords, API keys, tokens) ne doivent JAMAIS être:
\begin{itemize}
    \item Directement dans le code
    \item Commités dans version control
    \item Loggés ou printés
    \item Visibles en plain text
\end{itemize}

\subsection{Gestion de Secrets}

\subsubsection{Variables d'Environnement}

\lstset{language=bash}
\begin{lstlisting}
export DB_PASSWORD="secret_password"
export API_KEY="secret_apikey"
python app.py
\end{lstlisting}

\subsubsection{Kubernetes Secrets}

\lstset{language=yaml}
\begin{lstlisting}
apiVersion: v1
kind: Secret
metadata:
  name: app-secrets
type: Opaque
stringData:
  db_password: "secure_password"
  api_key: "secure_key"
---
apiVersion: v1
kind: Pod
metadata:
  name: app
spec:
  containers:
  - name: app
    image: app:latest
    env:
    - name: DB_PASSWORD
      valueFrom:
        secretKeyRef:
          name: app-secrets
          key: db_password
\end{lstlisting}

\subsubsection{Ansible Vault}

Chiffrer données sensibles en Ansible:

\lstset{language=bash}
\begin{lstlisting}
# Créer fichier vault
ansible-vault create secrets.yml

# Éditer fichier vault
ansible-vault edit secrets.yml

# Utiliser dans playbook
ansible-playbook playbook.yml --ask-vault-pass
\end{lstlisting}

\section{Bonnes Pratiques de Sécurité}

\begin{enumerate}
    \item \textbf{Least Privilege}: Chaque composant a permissions minimales nécessaires
    \item \textbf{Defense in Depth}: Multiples couches de sécurité
    \item \textbf{Regular Updates}: Maintenir systèmes à jour
    \item \textbf{Monitoring}: Audit trails et logging détaillé
    \item \textbf{Network Isolation}: Segmenter réseau avec firewalls
    \item \textbf{RBAC}: Role-Based Access Control
    \item \textbf{Encryption}: Données en transit et au repos
\end{enumerate}

\newpage

% CHAPITRE 10: Monitoring et Logging
\chapter{Monitoring, Logging et Observabilité}

\section{Importance du Monitoring}

\subsection{Types de Monitoring}

\begin{itemize}
    \item \textbf{Infrastructure}: CPU, Mémoire, Disque, Réseau
    \item \textbf{Application}: Temps réponse, Erreurs, Throughput
    \item \textbf{Business}: Utilisateurs actifs, Transactions
    \item \textbf{Security}: Tentatives accès non-autorisés
\end{itemize}

\section{Logging et Observabilité}

\subsection{Structured Logging}

Les logs doivent être structurés en JSON:

\lstset{language=python}
\begin{lstlisting}
import logging
import json
from datetime import datetime

class JSONFormatter(logging.Formatter):
    def format(self, record):
        log_data = {
            'timestamp': datetime.utcnow().isoformat(),
            'level': record.levelname,
            'logger': record.name,
            'message': record.getMessage(),
            'module': record.module,
            'function': record.funcName,
            'line': record.lineno
        }
        return json.dumps(log_data)

handler = logging.StreamHandler()
handler.setFormatter(JSONFormatter())
logger = logging.getLogger(__name__)
logger.addHandler(handler)

logger.info("Application started", extra={"version": "1.0.0"})
\end{lstlisting}

\subsection{Kubernetes Logging}

Récupérer logs:

\lstset{language=bash}
\begin{lstlisting}
# Logs d'un pod
kubectl logs deployment/devops-app

# Logs en temps réel
kubectl logs -f deployment/devops-app

# Logs dernières N lignes
kubectl logs deployment/devops-app --tail=100

# Logs d'un conteneur spécifique
kubectl logs deployment/devops-app -c app
\end{lstlisting}

\section{Stack ELK (Elasticsearch, Logstash, Kibana)}

Architecture de logging distribuée:

\lstset{language=yaml}
\begin{lstlisting}
version: '3.8'

services:
  elasticsearch:
    image: docker.elastic.co/elasticsearch/elasticsearch:7.14.0
    environment:
      - discovery.type=single-node
      - xpack.security.enabled=false
    ports:
      - "9200:9200"

  logstash:
    image: docker.elastic.co/logstash/logstash:7.14.0
    volumes:
      - ./logstash.conf:/usr/share/logstash/pipeline/logstash.conf
    ports:
      - "5000:5000"
    depends_on:
      - elasticsearch

  kibana:
    image: docker.elastic.co/kibana/kibana:7.14.0
    ports:
      - "5601:5601"
    depends_on:
      - elasticsearch
\end{lstlisting}

\newpage

% CHAPITRE 11: Documentation et Reproduction
\chapter{Documentation Complète et Reproductibilité}

\section{Guide de Installation et Déploiement}

\subsection{Prérequis}

\begin{itemize}
    \item Docker et Docker Compose (v20.10+)
    \item Terraform (v1.0+)
    \item Ansible (v2.10+)
    \item kubectl (v1.24+)
    \item Helm (v3.0+)
    \item Git
    \item Une clé AWS (ou autre cloud provider)
\end{itemize}

\subsection{Installation Locale avec Docker Compose}

\lstset{language=bash}
\begin{lstlisting}
# Cloner le repository
git clone https://github.com/votre-org/devops-project.git
cd devops-project

# Démarrer l'application localement
docker-compose up -d

# Vérifier santé application
curl http://localhost:5000/health

# Voir les logs
docker-compose logs -f app

# Arrêter l'application
docker-compose down
\end{lstlisting}

\subsection{Déploiement sur Kubernetes (Cloud)}

\lstset{language=bash}
\begin{lstlisting}
# 1. Initialiser Terraform
cd terraform
terraform init

# 2. Voir plan
terraform plan -out=tfplan

# 3. Appliquer infrastructure
terraform apply tfplan

# 4. Configurer kubectl
aws eks update-kubeconfig --region us-east-1 --name devops-app-cluster

# 5. Vérifier connexion
kubectl get nodes

# 6. Créer namespace
kubectl create namespace production

# 7. Déployer application
kubectl apply -f ../kubernetes/deployment.yaml
kubectl apply -f ../kubernetes/service.yaml

# 8. Vérifier déploiement
kubectl get deployment -n production
kubectl get pods -n production

# 9. Attendre disponibilité
kubectl rollout status deployment/devops-app -n production

# 10. Tester application
kubectl port-forward -n production svc/devops-app 8080:80
curl http://localhost:8080/health
\end{lstlisting}

\subsection{Configuration Ansible}

\lstset{language=bash}
\begin{lstlisting}
# Entrer répertoire Ansible
cd ansible

# Vérifier connectivité aux hôtes
ansible all -i inventory.ini -m ping

# Exécuter playbook complet
ansible-playbook -i inventory.ini playbook.yml

# Exécuter sur groupe spécifique
ansible-playbook -i inventory.ini playbook.yml -l app
\end{lstlisting}

\section{Structure de Repository}

\begin{lstlisting}
devops-project/
├── app/                          # Code application
│   ├── app.py
│   ├── requirements.txt
│   ├── templates/
│   │   └── index.html
│   └── static/
│       ├── style.css
│       └── script.js
├── kubernetes/                   # Manifests K8s
│   ├── deployment.yaml
│   └── service.yaml
├── terraform/                    # Infrastructure as Code
│   ├── main.tf
│   ├── providers.tf
│   ├── variables.tf
│   └── outputs.tf
├── ansible/                      # Configuration Management
│   ├── playbook.yml
│   └── inventory.ini
├── .github/
│   └── workflows/               # CI/CD Pipelines
│       └── ci-cd.yml
├── Dockerfile                   # Container image
├── docker-compose.yml           # Local dev setup
├── README.md                    # Guide principal
└── docs/                        # Documentation supplémentaire
    ├── ARCHITECTURE.md
    ├── DEPLOYMENT.md
    ├── TROUBLESHOOTING.md
    └── BEST_PRACTICES.md
\end{lstlisting}

\section{Checklist de Déploiement}

Avant tout déploiement production:

\begin{enumerate}
    \item[] \checkmark Infrastructure provisionnée et testée
    \item[] \checkmark Tous les tests passent (unitaires + intégration)
    \item[] \checkmark Image Docker build et push vers registry
    \item[] \checkmark Kubernetes manifests validés
    \item[] \checkmark Secrets gérés via sécurité appropriée
    \item[] \checkmark Monitoring et logging configurés
    \item[] \checkmark Backup et disaster recovery plan
    \item[] \checkmark Documentation mise à jour
    \item[] \checkmark Runbooks pour opérations
    \item[] \checkmark Communication aux stakeholders
\end{enumerate}

\newpage

% CHAPITRE 12: Troubleshooting et Maintenance
\chapter{Troubleshooting, Maintenance et Operations}

\section{Problèmes Communs et Solutions}

\subsection{Problèmes Docker}

\subsubsection{Conteneur ne démarre pas}

\lstset{language=bash}
\begin{lstlisting}
# Voir les logs
docker logs <container_id>

# Inspecter détails conteneur
docker inspect <container_id>

# Tester image en interactif
docker run -it <image> /bin/bash

# Vérifier ressources
docker stats
\end{lstlisting}

\subsubsection{Problèmes de connectivité}

\lstset{language=bash}
\begin{lstlisting}
# Vérifier réseau
docker network ls
docker network inspect <network>

# Tester connectivité entre conteneurs
docker exec <container> ping <autre_container>
\end{lstlisting}

\subsection{Problèmes Kubernetes}

\subsubsection{Pod ne démarre pas}

\lstset{language=bash}
\begin{lstlisting}
# Voir statut pod
kubectl describe pod <pod_name> -n <namespace>

# Voir logs
kubectl logs <pod_name> -n <namespace>

# Voir événements
kubectl get events -n <namespace> --sort-by='.lastTimestamp'

# Utiliser temporaire pod pour debug
kubectl run debug --image=busybox -it -- sh
\end{lstlisting}

\subsubsection{Service inaccessible}

\lstset{language=bash}
\begin{lstlisting}
# Vérifier service
kubectl get svc -n <namespace>
kubectl describe svc <service_name> -n <namespace>

# Vérifier endpoints
kubectl get endpoints -n <namespace>

# Test depuis pod
kubectl run curl --image=curlimages/curl -it -- sh
curl http://<service_name>:<port>

# Port-forward debug
kubectl port-forward service/<service_name> 8080:80
\end{lstlisting}

\subsection{Problèmes Terraform}

\subsubsection{Drift de state}

\lstset{language=bash}
\begin{lstlisting}
# Voir le state actuel
terraform show

# Rafraîchir state
terraform refresh

# Importer ressource manquante
terraform import aws_instance.example i-1234567890abcdef0

# Nettoyer ressource du state
terraform state rm aws_instance.example
\end{lstlisting}

\section{Maintenance Régulière}

\subsection{Updates du Système}

\begin{itemize}
    \item Scheduler updates des base images mensuels
    \item Vérifier vulnerabilités avec tools comme Trivy
    \item Mettre à jour dépendances application
    \item Auditer accès secrets régulièrement
\end{itemize}

\subsection{Backups et Disaster Recovery}

\begin{itemize}
    \item Backup Terraform state en S3 avec versioning
    \item Backup données applicatives via snapshots
    \item Tester restaurations régulièrement
    \item Documenter procedures de recovery
\end{itemize}

\section{Operational Runbooks}

\subsection{Scaler Application}

\lstset{language=bash}
\begin{lstlisting}
# Scaler replicas Deployment
kubectl scale deployment devops-app --replicas=5 -n production

# Vérifier scaling
kubectl get deployment devops-app -n production
kubectl get pods -n production
\end{lstlisting}

\subsection{Rollback Déploiement}

\lstset{language=bash}
\begin{lstlisting}
# Voir historique rollout
kubectl rollout history deployment/devops-app -n production

# Rollback à version précédente
kubectl rollout undo deployment/devops-app -n production

# Rollback à version spécifique
kubectl rollout undo deployment/devops-app --to-revision=2 -n production
\end{lstlisting}

\section{Monitoring Santé Infrastructure}

\subsection{Health Checks}

\lstset{language=bash}
\begin{lstlisting}
#!/bin/bash
# Scénarios de test basiques

echo "Testing Docker..."
docker ps > /dev/null && echo "✓ Docker OK" || echo "✗ Docker FAILED"

echo "Testing Kubernetes..."
kubectl cluster-info > /dev/null && echo "✓ Kubernetes OK" || echo "✗ Kubernetes FAILED"

echo "Testing Application..."
APP_HEALTH=$(curl -s http://app-service/health | jq .status)
[ "$APP_HEALTH" == "healthy" ] && echo "✓ Application OK" || echo "✗ Application FAILED"

echo "Testing Database..."
psql -h db-host -U user -d database -c "SELECT 1" > /dev/null && echo "✓ Database OK" || echo "✗ Database FAILED"
\end{lstlisting}

\newpage

% CHAPITRE 13: Conclusions et Recommandations
\chapter{Conclusions et Recommandations}

\section{Résumé des Réalisations}

Ce projet a démontré la mise en place réussie d'une infrastructure complète DevOps incluant:

\begin{enumerate}
    \item Application web containerisée et scalable
    \item Infrastructure cloud reproductible via Terraform
    \item Automatisation complète des déploiements
    \item Orchestration Kubernetes pour haute disponibilité
    \item Pipeline CI/CD entièrement automatisée
    \item Documentation exhaustive et guides opérationnels
\end{enumerate}

\section{Points Clés de la Solution}

\subsection{Architecture Moderne}

L'architecture suit les principes cloud-native:
\begin{itemize}
    \item Microservices et conteneurs
    \item Stateless applications
    \item Infrastructure as Code
    \item Immutable infrastructure
    \item Decentralized governance
\end{itemize}

\subsection{Automatisation Maximale}

Presque tout est automatisé:
\begin{itemize}
    \item Provisioning infrastructure
    \item Configuration servers
    \item Build et tests
    \item Déploiement application
    \item Scalabilité
    \item Monitoring et alertes
\end{itemize}

\section{Recommandations Futures}

\subsection{Court Terme (1-3 mois)}

\begin{enumerate}
    \item Implémenter monitoring avec Prometheus et Grafana
    \item Ajouter centralized logging avec ELK
    \item Configurer alerting via PagerDuty/Opsgenie
    \item Implémenter security scanning en CI/CD
    \item Ajouter tests de charge et performance
\end{enumerate}

\subsection{Moyen Terme (3-6 mois)}

\begin{enumerate}
    \item Implémenter service mesh (Istio) pour resilience
    \item Ajouter policy management (OPA/Gatekeeper)
    \item Mettre en place disaster recovery testing
    \item Implémenter A/B testing framework
    \item Ajouter cost optimization monitoring
\end{enumerate}

\subsection{Long Terme (6+ mois)}

\begin{enumerate}
    \item Expansion multi-region pour haute disponibilité globale
    \item Machine Learning ops (MLOps) pipelines
    \item Advanced cost allocation et chargeback
    \item Multi-cloud strategy pour vendor-lock-in reduction
    \item Serverless components où approprié
\end{enumerate}

\section{Lessons Learned}

\begin{itemize}
    \item \textbf{Infrastructure as Code importante}: Permet reproducibilité
    \item \textbf{Automatisation critical}: Erreurs humaines minimisées
    \item \textbf{Testing multiples layers}: Unitaire, intégration, e2e nécessaires
    \item \textbf{Monitoring from day 1}: Invaluable pour troubleshooting
    \item \textbf{Documentation vital}: Onboarding nouveau personnel facilité
    \item \textbf{Security baseline}: Ne pas ajouter en dernier
    \item \textbf{Gradual adoption}: Mieux que big-bang approach
\end{itemize}

\section{Conclusion Finale}

Ce projet démontre comment les technologies modernes peuvent être combinées pour créer une infrastructure robuste, scalable et maintenable. Les pratiques DevOps et Cloud-Native implémentées ici établissent une foundation solide pour:

\begin{itemize}
    \item Accélérer time-to-market
    \item Améliorer reliability et stability
    \item Réduire operational overhead
    \item Faciliter collaboration Dev-Ops
    \item Supporter croissance future
\end{itemize}

La clé du succès réside dans l'automatisation complète, la documentation claire, et l'adoption continue de meilleures pratiques au fur et à mesure que la technologie évolue.

\newpage

% Appendices
\chapter*{Appendix A: Commandes Référence}
\addcontentsline{toc}{chapter}{Appendix A: Commandes Référence}

\input{./append_commands.tex}

\chapter*{Appendix B: Configurations Complètes}
\addcontentsline{toc}{chapter}{Appendix B: Configurations Complètes}

Voir fichiers associés dans le repository pour configurations complètes.

\chapter*{Appendix C: Ressources Additionnelles}
\addcontentsline{toc}{chapter}{Appendix C: Ressources Additionnelles}

Ressources recommandées:
\begin{itemize}
    \item Kubernetes Official Documentation: https://kubernetes.io/docs/
    \item Terraform Registry: https://registry.terraform.io/
    \item Ansible Documentation: https://docs.ansible.com/
    \item Docker Best Practices: https://docs.docker.com/develop/
    \item Linux Academy et courses en ligne
\end{itemize}

\end{document}
